\chapter{Quantum Spin}

\section{The Stern-Gerlach experiment}
The Stern-Gerlach experiment (1922) demonstrated that a beam of silver atoms, despite having zero orbital angular momentum, splits into two distinct beams when passing through a non-uniform magnetic field $\vec{B}$. Silver atoms have this property because, except for a single external electron, their electronic configuration has spherical symmetry with $\vec{L}_{\text{TOT}}=0$. The external electron is in the $5s$ state ($\ell=0$), also contributing zero orbital angular momentum.


\begin{equation}
H=\frac{\vec{p}^{2}}{2 m}-\vec{\mu} \cdot \vec{B}
\end{equation}

Where $\vec{\mu}=-\frac{e}{2 m_{e} c} \vec{L}$ represents the magnetic moment of a hydrogen-like atom with a single electron. Silver atoms closely match this model. The unexpected splitting observed in the experiment arises from the force due to the magnetic moment interacting with a non-uniform magnetic field:

\begin{equation}
\vec{F}=-\vec{\nabla}(-\vec{\mu} \cdot \vec{B})=\mu_{3} \frac{\partial B_{3}}{\partial x_{3}} \hat{u}_{3}
\end{equation}

According to conventional quantum theory, since the outer electron has $\ell=0$, its magnetic moment should be:

\begin{equation}
\mu_{3}=-\frac{e}{2 m_{e} c} L_{3} \propto \hbar m=0
\end{equation}

With $m=0$ being the only possible eigenvalue for $L_3$ in s-orbital states. The observed splitting revealed that electrons possess an intrinsic angular momentum called spin $\vec{S}$, which contributes to the total magnetic moment. This spin has two possible states corresponding to eigenvalues $\pm\hbar/2$ of $S_3$, explaining the two-beam splitting pattern.

Numerous other experiments, including the anomalous Zeeman effect and helium spectroscopy, confirmed that spin is fundamental to atomic structure. This led to the spin-statistics theorem, which classifies elementary particles into fermions (half-integer spin particles like electrons, protons, and neutrons) and bosons (integer spin particles like photons or composite particles with even numbers of fermions).

\section{Mathematical formulation of spin}
While spin operators $S$ and angular momentum operators $L$ share similar algebraic properties, they differ fundamentally in their definitions:

\begin{itemize}
  \item $L$ is expressed through differential operators
  \item $S$ is defined using matrices
\end{itemize}

The spin operators are constructed from the Pauli matrices:

\[
\sigma_{1}=\left[\begin{array}{ll}
0 & 1  \\
1 & 0
\end{array}\right], \quad \sigma_{2}=\left[\begin{array}{cc}
0 & -i \\
i & 0
\end{array}\right], \quad \sigma_{3}=\left[\begin{array}{cc}
1 & 0 \\
0 & -1
\end{array}\right]
\]

The spin operators are defined as:

\begin{align}
S_{1} & =\frac{\hbar}{2} \sigma_{1} \\
S_{2} & =\frac{\hbar}{2} \sigma_{2}  \\
S_{3} & =\frac{\hbar}{2} \sigma_{3}
\end{align}

The Pauli matrices have the property:

\begin{equation}
\sigma_{1}^{2}=\mathbb{I}, \quad \sigma_{2}^{2}=\mathbb{I}, \quad \sigma_{3}^{2}=\mathbb{I}
\end{equation}

We can define raising and lowering operators for spin:

\begin{align}
& S_{+}=S_{1}+i S_{2}=\hbar\left[\begin{array}{ll}
0 & 1 \\
0 & 0
\end{array}\right] \\
& S_{-}=S_{1}-i S_{2}=\hbar\left[\begin{array}{ll}
0 & 0 \\
1 & 0
\end{array}\right]
\end{align}

Examining their commutation relations:

\begin{align}
S_1 S_2 - S_2 S_1 &= \frac{\hbar}{2}\left[\begin{array}{ll}
0 & 1 \\
1 & 0
\end{array}\right] \frac{\hbar}{2}\left[\begin{array}{cc}
0 & -i \\
i & 0
\end{array}\right] - \frac{\hbar}{2}\left[\begin{array}{cc}
0 & -i \\
i & 0
\end{array}\right] \frac{\hbar}{2}\left[\begin{array}{ll}
0 & 1 \\
1 & 0
\end{array}\right] \\
&= i\frac{\hbar^{2}}{4}\left[\begin{array}{cc}
1 & 0 \\
0 & -1
\end{array}\right] = i\hbar S_3
\end{align}


Indeed, the commutation relations for spin operators mirror those of angular momentum:

\begin{equation}
\left[S_{1}, S_{2}\right]=i \hbar S_{3}
\end{equation}

And by cyclic permutation:

\begin{align}
& {\left[S_{2}, S_{3}\right]=i \hbar S_{1}} \\
& {\left[S_{3}, S_{1}\right]=i \hbar S_{2}} \\
& {\left[S_{+}, S_{-}\right]=2 \hbar S_{3}}  \\
& {\left[S_{3}, S_{\pm}\right]=\pm \hbar S_{\pm}}
\end{align}

The total spin squared operator is:
\begin{equation}
S^{2}=S_{1}^{2}+S_{2}^{2}+S_{3}^{2}=\frac{3}{4} \hbar^{2} \mathbb{I}
\end{equation}

The standard basis consists of eigenstates of $S_3$, satisfying:

\begin{equation}
S_{3}\left|s, m_{s}\right\rangle=\hbar m_{s}\left|s, m_{s}\right\rangle \quad \text{with} \quad m_{s}=\pm\frac{1}{2}
\end{equation}

These eigenstates have the explicit matrix representation:

\begin{align}
& \left|s,+\frac{1}{2}\right\rangle=\left[\begin{array}{l}
1 \\
0
\end{array}\right]  \\
& \left|s,-\frac{1}{2}\right\rangle=\left[\begin{array}{l}
0 \\
1
\end{array}\right]
\end{align}

The eigenvalue equation for $S^2$ is:

\begin{equation}
S^{2}\left|s, m_{s}\right\rangle=\hbar^{2} s(s+1)\left|s, m_{s}\right\rangle
\end{equation}

From our matrix representation, we know:

\begin{equation}
S^{2}\left|s, m_{s}\right\rangle=\frac{3}{4} \hbar^{2} \mathbb{I}\left|s, m_{s}\right\rangle=\frac{3}{4} \hbar^{2}\left|s, m_{s}\right\rangle
\end{equation}

This gives us the value of $s$:

\begin{equation}
s(s+1)=\frac{3}{4} \Longrightarrow s=\frac{1}{2}
\end{equation}

We select the positive solution by analogy with angular momentum. Since $s$ has only one value, we can simplify the notation:

\begin{equation}
|s, \pm 1/2\rangle=|\pm 1/2\rangle
\end{equation}

The action of the raising and lowering operators on these basis states is:

\[
S_{+}|-1/2\rangle=\hbar\left[\begin{array}{ll}
0 & 1  \\
0 & 0
\end{array}\right]\left[\begin{array}{l}
0 \\
1
\end{array}\right]=\hbar\left[\begin{array}{l}
1 \\
0
\end{array}\right]=\hbar|+1/2\rangle
\]

And similarly:

\[
S_{-}|+1/2\rangle=\hbar\left[\begin{array}{ll}
0 & 0  \\
1 & 0
\end{array}\right]\left[\begin{array}{l}
1 \\
0
\end{array}\right]=\hbar\left[\begin{array}{l}
0 \\
1
\end{array}\right]=\hbar|-1/2\rangle
\]


The remaining operations yield zero results:
\begin{align}
S_+|+1/2\rangle = 0\\
S_-|-1/2\rangle = 0
\end{align}

When dealing with spin-dependent Hamiltonians, we work with spinor wavefunctions:

\[
|\Psi\rangle=\sum_{m_{s}= \pm 1/2} \psi_{m_{s}}\left|m_{s}\right\rangle=\psi_{+}|+1/2\rangle+\psi_{-}|-1/2\rangle=\left[\begin{array}{l}
\psi_{+}  \\
\psi_{-}
\end{array}\right]
\]

For a Hamiltonian of the form:

\begin{equation}
\mathcal{H}=\mathcal{H}_{0}+\vec{\alpha} \cdot(\vec{L}+g \vec{S})
\end{equation}

The corresponding quantum mechanical operator is:

\begin{equation}
H=H_{0}+\vec{\alpha} \cdot(\vec{L}+g \vec{S})
\end{equation}

Where:

\[
\vec{L}=\left[\begin{array}{l}
L_{1}  \\
L_{2} \\
L_{3}
\end{array}\right], \quad \vec{S}=\left[\begin{array}{l}
S_{1} \\
S_{2} \\
S_{3}
\end{array}\right]
\]

Here $\vec{L}$ and $\vec{S}$ represent "vectors of operators" - a notation for operations like scalar products where each component operator acts on the corresponding component of another vector.

Applying this Hamiltonian to a spinor state:

\begin{align}
& \left(H_{0}+\vec{\alpha} \cdot(\vec{L}+g \vec{S})\right)\left[\begin{array}{l}
\psi_{+} \\
\psi_{-}
\end{array}\right]= \\
& =\left[\begin{array}{l}
H_{0} \psi_{+} \\
H_{0} \psi_{-}
\end{array}\right]+\left[\begin{array}{l}
\vec{\alpha} \cdot \vec{L} \psi_{+} \\
\vec{\alpha} \cdot \vec{L} \psi_{-}
\end{array}\right]+\frac{g \hbar}{2}\left[\begin{array}{cc}
\alpha_{3} & \alpha_{1}-i \alpha_{2} \\
\alpha_{1}+i \alpha_{2} & -\alpha_{3}
\end{array}\right]\left[\begin{array}{l}
\psi_{+} \\
\psi_{-}
\end{array}\right]=  \\
& =\left[\begin{array}{l}
H_{0} \psi_{+} \\
H_{0} \psi_{-}
\end{array}\right]+\left[\begin{array}{l}
\vec{\alpha} \cdot \vec{L} \psi_{+} \\
\vec{\alpha} \cdot \vec{L} \psi_{-}
\end{array}\right]+\frac{g \hbar}{2}\left[\begin{array}{l}
\alpha_{3} \psi_{+}+\left(\alpha_{1}-i \alpha_{2}\right) \psi_{-} \\
\left(\alpha_{1}+i \alpha_{2}\right) \psi_{+}-\alpha_{3} \psi_{-}
\end{array}\right]
\end{align}

This leads to the coupled Schrödinger equations for the spinor components:

\[
-i \hbar \frac{\partial}{\partial t}\left[\begin{array}{l}
\psi_{+}  \\
\psi_{-}
\end{array}\right]=\left[\begin{array}{l}
\left(H_{0}+\vec{\alpha} \cdot \vec{L}+\frac{g\hbar}{2}\alpha_{3}\right) \psi_{+}+\frac{g\hbar}{2}\left(\alpha_{1}-i \alpha_{2}\right) \psi_{-} \\
\left(H_{0}+\vec{\alpha} \cdot \vec{L}-\frac{g\hbar}{2}\alpha_{3}\right) \psi_{-}+\frac{g\hbar}{2}\left(\alpha_{1}+i \alpha_{2}\right) \psi_{+}
\end{array}\right]
\]


These equations form a coupled system of differential equations.

Setting $\vec{\alpha}=e \vec{B} / 2 m_{e} c$, we obtain the Pauli equation, which represents the non-relativistic limit of the Dirac equation:

\[
-i \hbar \frac{\partial}{\partial t}\left[\begin{array}{l}
\psi_{+}  \\
\psi_{-}
\end{array}\right]=\left(-\frac{\hbar^{2}}{2 m_{e}} \nabla^{2}+V(\vec{x})+\frac{e}{2 m_{e} c} \vec{B} \cdot(\vec{L}+g \vec{S})\right)\left[\begin{array}{l}
\psi_{+} \\
\psi_{-}
\end{array}\right]
\]

\section{Spin-1/2 free particle in a magnetic field}
For a point-like particle with spin 1/2 and mass $m$ in a magnetic field, the Hamiltonian simplifies to:

\begin{equation}
H=H_{k}+H_{s}=\frac{p^{2}}{2 m}+k \vec{B} \cdot \vec{S}
\end{equation}

Using spherical coordinates for the magnetic field:

\begin{equation}
B_{1}=B \sin \theta \cos \phi, \quad B_{2}=B \sin \theta \sin \phi, \quad B_{3}=B \cos \theta
\end{equation}

The spin-dependent term becomes:

\begin{align}
k \vec{B} \cdot \vec{S} &= k \frac{\hbar}{2}\left[B_{1}\left(\begin{array}{ll}
0 & 1 \\
1 & 0
\end{array}\right)+B_{2}\left(\begin{array}{cc}
0 & -i \\
i & 0
\end{array}\right)+B_{3}\left(\begin{array}{cc}
1 & 0 \\
0 & -1
\end{array}\right)\right] \\
&= \frac{\hbar}{2} k B\left[\begin{array}{cc}
\cos \theta & e^{-i \phi} \sin \theta \\
e^{i \phi} \sin \theta & -\cos \theta
\end{array}\right]
\end{align}

This model describes a neutral particle (like a neutron) in a time-dependent magnetic field $\vec{B}$ where $\theta=\theta(t)$ and $\phi=\phi(t)$. Such systems exhibit the Berry phase - an additional phase in the wavefunction arising from the time-dependence of $\vec{B}$.

The energy eigenvalue problem is:

\[
\left[\frac{p^{2}}{2 m}+k \vec{B} \cdot \vec{S}\right]\left[\begin{array}{l}
\psi_{+}  \\
\psi_{-}
\end{array}\right]=E\left[\begin{array}{l}
\psi_{+} \\
\psi_{-}
\end{array}\right]
\]

The spinor wavefunction can be factorized into a plane wave and a spin vector:

\[
|\psi\rangle=\left[\begin{array}{l}
\psi_{+}  \\
\psi_{-}
\end{array}\right]=e^{i \vec{k} \cdot \vec{x}}|v\rangle=e^{i \vec{k} \cdot \vec{x}}\left[\begin{array}{l}
v_{+} \\
v_{-}
\end{array}\right], \quad \psi_{\pm}=v_{\pm} e^{i \vec{k} \cdot \vec{x}}
\]

Since the kinetic energy operator commutes with the spin operators, the eigenvectors of $k \vec{B} \cdot \vec{S}$ are:

\[
|v(+1)\rangle=\left[\begin{array}{c}
\cos (\theta/2) e^{-i \phi/2}  \\
\sin (\theta/2) e^{i \phi/2}
\end{array}\right], \quad|v(-1)\rangle=\left[\begin{array}{c}
-\sin (\theta/2) e^{-i \phi/2} \\
\cos (\theta/2) e^{i \phi/2}
\end{array}\right]
\]

These satisfy the eigenvalue equation:

\[
\frac{\hbar}{2} k B\left[\begin{array}{cc}
\cos \theta & e^{-i \phi} \sin \theta  \\
e^{i \phi} \sin \theta & -\cos \theta
\end{array}\right]|v(\mu)\rangle=\mu \frac{\hbar}{2} k B|v(\mu)\rangle
\]

where $\mu = \pm 1$ corresponds to spin aligned or anti-aligned with the magnetic field.


From these eigenvectors, we can determine the energy eigenvalues:

\begin{equation}
\left[\frac{p^{2}}{2 m}+k \vec{B} \cdot \vec{S}\right] e^{i \vec{k} \cdot \vec{x}}|v(\mu)\rangle=|v(\mu)\rangle \frac{p^{2}}{2 m} e^{i \vec{k} \cdot \vec{x}}+e^{i \vec{k} \cdot \vec{x}} k \vec{B} \cdot \vec{S}|v(\mu)\rangle=\left(\frac{\hbar^{2} k^{2}}{2 m}+\mu \frac{\hbar}{2} k B\right) e^{i \vec{k} \cdot \vec{x}}|v(\mu)\rangle
\end{equation}

Thus, the energy spectrum is:

\begin{equation}
E(\vec{k}, \mu)=\frac{\hbar^{2} k^{2}}{2 m}+\mu \frac{\hbar}{2} k B
\end{equation}

This shows how the energy splits into two levels due to the magnetic field, with the separation proportional to the field strength.

\section{Addition of angular momenta}
When two angular momenta interact in a quantum system, we need to understand how to combine them. Classically, we can add vectors $\vec{J}=\vec{L}+\vec{S}$ or $\vec{J}=\vec{S}_1+\vec{S}_2$. In quantum mechanics, this corresponds to the operator $\vec{J}=\vec{L}+\vec{S}$, which satisfies:

\begin{align}
\left[J_{n}, J_{m}\right] &= \left[L_{n}+S_{n}, L_{m}+S_{m}\right]\\
&= \left[L_{n}, L_{m}\right]+\left[L_{n}, S_{m}\right]+\left[S_{n}, L_{m}\right]+\left[S_{n}, S_{m}\right]\\
&= i \hbar \epsilon_{n m k} L_{k}+i \hbar \epsilon_{n m k} S_{k}=i \hbar \epsilon_{n m k} J_{k}
\end{align}

This shows that $\vec{J}$ satisfies the same commutation relations as $\vec{L}$ and $\vec{S}$, confirming that $\left[J^{2}, J_{3}\right]=0$. Therefore, the quantum numbers $j$ and $m_j$ provide a good basis.

Classically, the magnitudes of these vectors are:

\[
\begin{array}{r}
|\vec{L}|=\hbar \ell \\
|\vec{S}|=\hbar s  \\
\hbar(\ell-s) \leq|\vec{J}| \leq \hbar(\ell+s)
\end{array}
\]

In quantum mechanics, the operators have eigenvalues $\hbar^{2}\ell(\ell+1)$ and $\hbar^{2}s(s+1)$, giving "quantum lengths":

\begin{align}
|L| &= \sqrt{\hbar^{2} \ell(\ell+1)} \sim \hbar \ell  \\
|S| &= \sqrt{\hbar^{2} s(s+1)} \sim \hbar s
\end{align}

The quantum number $j$ must satisfy:

\begin{equation}
-j \leq m_{j} \leq j, \quad \ell-s \leq j \leq \ell+s
\end{equation}

Examining the commutators between $J^2$ and the individual components $L_3$ or $S_3$:

\begin{align}
\left[J^{2}, S_{3}\right] &= \left[(L+S)^{2}, S_{3}\right]\\
&= \left[L^{2}+S^{2}+2 \vec{L} \cdot \vec{S}, S_{3}\right]\\
&= \left[L^{2}, S_{3}\right]+\left[S^{2}, S_{3}\right]+2\left[\vec{L} \cdot \vec{S}, S_{3}\right]\\
&= 2 \sum_{i}\left[L_{i} S_{i}, S_{3}\right]\\
&= 2 \sum_{i}\left(L_{i}\left[S_{i}, S_{3}\right]+\left[L_{i}, S_{3}\right] S_{i}\right)\\
&= 2 \sum_{i} L_{i}\left[S_{i}, S_{3}\right] \neq 0
\end{align}

Similarly:

\begin{align}
\left[J^{2}, L_{3}\right] &= \left[(L+S)^{2}, L_{3}\right]\\
&= 2 \sum_{i}\left[L_{i}, L_{3}\right] S_{i} \neq 0
\end{align}

However, the sum of these commutators is zero:

\begin{align}
&2 \sum_{i}\left[L_{i}, L_{3}\right] S_{i}+2 \sum_{i} L_{i}\left[S_{i}, S_{3}\right]\\
&= 2\left(\left[L_{1}, L_{3}\right] S_{1}+\left[L_{2}, L_{3}\right] S_{2}+L_{1}\left[S_{1}, S_{3}\right]+L_{2}\left[S_{2}, S_{3}\right]\right)\\
&= 2\left(-i \hbar L_{2} S_{1}+i \hbar L_{1} S_{2}-i \hbar L_{1} S_{2}+i \hbar L_{2} S_{1}\right)=0
\end{align}

This confirms that $\left[J^{2}, J_{k}\right]=0$. Other useful commutators include:

\begin{align}
&\left[J^{2}, L^{2}\right]=0 \\
&\left[L^{2}, J_{3}\right]=0  \\
&\left[J^{2}, S^{2}\right]=0 \\
&\left[S^{2}, J_{3}\right]=0
\end{align}

These commutation relations allow us to define a complete set of commuting observables:

\begin{equation}
S^{2}, L^{2}, J^{2}, J_{3} \quad \Longleftrightarrow \quad s, \ell, j, m_{j}
\end{equation}

The states with $m_j = m + m_s$ can be formed from different combinations of $m$ and $m_s$ that sum to the same value. This degeneracy means that states in the coupled basis $|j, m_j, \ell, s\rangle$ are superpositions of the uncoupled basis states:

\begin{equation}
\left|j, m_{j}, \ell, s\right\rangle=\sum_{m, m_{s}}^{*} C\left(m, m_{s}\right)|\ell, m\rangle\left|s, m_{s}\right\rangle
\end{equation}


The symbol * implies that we keep $m+m_{s}$ constant. This state satisfies:

\begin{align}
L^{2}\left|j, m_{j}, \ell, s\right\rangle &= \sum_{m, m_{s}}^{*} C\left(m, m_{s}\right) L^{2}|\ell, m\rangle\left|s, m_{s}\right\rangle \\
&= \sum_{m, m_{s}}^{*} C\left(m, m_{s}\right)\left(\hbar^{2} \ell(\ell+1)\right)|\ell, m\rangle\left|s, m_{s}\right\rangle \\
&= \hbar^{2} \ell(\ell+1) \sum_{m, m_{s}}^{*} C\left(m, m_{s}\right)|\ell, m\rangle\left|s, m_{s}\right\rangle=\hbar^{2} \ell(\ell+1)\left|j, m_{j}, \ell, s\right\rangle
\end{align}

Similarly for the spin operator:

\begin{align}
S^{2}\left|j, m_{j}, \ell, s\right\rangle &= \sum_{m, m_{s}}^{*} C\left(m, m_{s}\right)|\ell, m\rangle S^{2}\left|s, m_{s}\right\rangle \\
&= \sum_{m, m_{s}}^{*} C\left(m, m_{s}\right)\left(\hbar^{2} s(s+1)\right)|\ell, m\rangle\left|s, m_{s}\right\rangle \\
&= \hbar^{2} s(s+1) \sum_{m, m_{s}}^{*} C\left(m, m_{s}\right)|\ell, m\rangle\left|s, m_{s}\right\rangle=\hbar^{2} s(s+1)\left|j, m_{j}, \ell, s\right\rangle
\end{align}

For the z-component of total angular momentum:

\begin{align}
J_{3}\left|j, m_{j}, \ell, s\right\rangle &= \sum_{m, m_{s}}^{*} C\left(m, m_{s}\right)\left(L_{3}|\ell, m\rangle\left|s, m_{s}\right\rangle+|\ell, m\rangle S_{3}\left|s, m_{s}\right\rangle\right) \\
&= \sum_{m, m_{s}}^{*} C\left(m, m_{s}\right) \hbar\left(m+m_{s}\right)|\ell, m\rangle\left|s, m_{s}\right\rangle \\
&= \hbar m_{j} \sum_{m, m_{s}}^{*} C\left(m, m_{s}\right)|\ell, m\rangle\left|s, m_{s}\right\rangle=\hbar m_{j}\left|j, m_{j}, \ell, s\right\rangle
\end{align}

The coupled basis states for $j=\ell\pm\frac{1}{2}$ can be expressed as:

\[
\left\{\begin{array}{l}
\left|\ell-\frac{1}{2}, m_{j}, \ell, s\right\rangle=\alpha_{-}\left|\ell, m_{j}-\frac{1}{2}\right\rangle\left|s,+\frac{1}{2}\right\rangle+\beta_{-}\left|\ell, m_{j}+\frac{1}{2}\right\rangle\left|s,-\frac{1}{2}\right\rangle  \\
\left|\ell+\frac{1}{2}, m_{j}, \ell, s\right\rangle=\alpha_{+}\left|\ell, m_{j}-\frac{1}{2}\right\rangle\left|s,+\frac{1}{2}\right\rangle+\beta_{+}\left|\ell, m_{j}+\frac{1}{2}\right\rangle\left|s,-\frac{1}{2}\right\rangle
\end{array}\right.
\]

The coefficients are given by:

\begin{equation}
\alpha_{\pm}= \pm \sqrt{\frac{\ell \pm m_{j}+\frac{1}{2}}{2\ell+1}}= \pm \beta_{\mp}
\end{equation}

These coefficients satisfy the normalization condition $\left|\alpha_{\pm}\right|^{2}+\left|\beta_{\pm}\right|^{2}=1$. For instance:

\begin{align}
\alpha_{+} &= \sqrt{\frac{\ell+m_{j}+\frac{1}{2}}{2\ell+1}}  \\
\beta_{+} &= -\sqrt{\frac{\ell-m_{j}+\frac{1}{2}}{2\ell+1}}
\end{align}

Verifying normalization:

\begin{equation}
\left|\alpha_{+}\right|^{2}+\left|\beta_{+}\right|^{2}=\frac{\ell+m_{j}+\frac{1}{2}}{2\ell+1}+\frac{\ell-m_{j}+\frac{1}{2}}{2\ell+1}=\frac{2\ell+1}{2\ell+1}=1
\end{equation}

\section{Coupling arbitrary angular momenta}
For two general angular momentum operators $J_{a}$ and $J_{b}$ with eigenstates:

\begin{align}
J_{a3}\left|j_{a}, m_{a}\right\rangle &= \hbar m_{a}\left|j_{a}, m_{a}\right\rangle \\
J_{b3}\left|j_{b}, m_{b}\right\rangle &= \hbar m_{b}\left|j_{b}, m_{b}\right\rangle \\
J_{a}^{2}\left|j_{a}, m_{a}\right\rangle &= \hbar^2 j_{a}\left(j_{a}+1\right)\left|j_{a}, m_{a}\right\rangle  \\
J_{b}^{2}\left|j_{b}, m_{b}\right\rangle &= \hbar^{2} j_{b}\left(j_{b}+1\right)\left|j_{b}, m_{b}\right\rangle
\end{align}

The coupled states are formed as:

\begin{equation}
\left|j, m_{j}, j_{a}, j_{b}\right\rangle=\sum_{m_{a}, m_{b}}^{*} C\left(m_{a}, m_{b}\right)\left|j_{a}, m_{a}\right\rangle\left|j_{b}, m_{b}\right\rangle
\end{equation}

These states diagonalize the total angular momentum operators:

\begin{align}
J_{3}\left|j, m_{j}, j_{a}, j_{b}\right\rangle &= \hbar m_{j}\left|j, m_{j}, j_{a}, j_{b}\right\rangle \\
J^{2}\left|j, m_{j}, j_{a}, j_{b}\right\rangle &= \hbar^2 j(j+1)\left|j, m_{j}, j_{a}, j_{b}\right\rangle
\end{align}

\section{Spin coupling illustration}
As an illustrative example, consider coupling two spin-$\frac{1}{2}$ systems $S_{a}$ and $S_{b}$. Here $j_{a}=j_{b}=\frac{1}{2}$, so the possible values of total angular momentum $j$ are:

\begin{equation}
j=0,1
\end{equation}

And the corresponding $m_j$ values are:

\begin{equation}
m_{j}=-1,0,1
\end{equation}


\section{Analysis of coupled spin states}
For the case $j=0$, we have $m_j=0$. This state can be formed from two combinations: $m_a=\frac{1}{2}, m_b=-\frac{1}{2}$ or $m_a=-\frac{1}{2}, m_b=\frac{1}{2}$. This configuration yields the singlet state:

\begin{equation}
|0,0,\frac{1}{2},\frac{1}{2}\rangle=\frac{1}{\sqrt{2}}(|\frac{1}{2},-\frac{1}{2}\rangle|\frac{1}{2},+\frac{1}{2}\rangle+\sigma|\frac{1}{2},+\frac{1}{2}\rangle|\frac{1}{2},-\frac{1}{2}\rangle)
\end{equation}

While both values $\sigma=\pm 1$ satisfy normalization, we must determine which gives the correct eigenvalue for $J^2$. We require $J^2|0,0,\frac{1}{2},\frac{1}{2}\rangle=0$. First, we express $J^2$ in terms of ladder operators:

\begin{equation}
J^{2}=J_{3}^{2}+J_{1}^{2}+J_{2}^{2}=J_{3}^{2}+\frac{1}{2}(J_{+}J_{-}+J_{-}J_{+})=J_{3}^{2}+\frac{1}{2}(2J_{+}J_{-}-2\hbar J_{3})
\end{equation}

Since $J_3|0,0,\frac{1}{2},\frac{1}{2}\rangle=0$, the terms involving $J_3$ vanish, leaving:

\begin{equation}
J^{2}=J_{+}J_{-}=(S_{a+}+S_{b+})(S_{a-}+S_{b-})
\end{equation}

Applying this to our state:

\begin{align}
J^{2}|0,0,\frac{1}{2},\frac{1}{2}\rangle &= (S_{a+}+S_{b+})(S_{a-}+S_{b-}) \frac{1}{\sqrt{2}}(|-\frac{1}{2}\rangle|+\frac{1}{2}\rangle+\sigma|+\frac{1}{2}\rangle|-\frac{1}{2}\rangle) \\
&= (S_{a+}+S_{b+}) \frac{1}{\sqrt{2}}(S_{a-}|-\frac{1}{2}\rangle|+\frac{1}{2}\rangle+\sigma S_{a-}|+\frac{1}{2}\rangle|-\frac{1}{2}\rangle+|-\frac{1}{2}\rangle S_{b-}|+\frac{1}{2}\rangle+\sigma|+\frac{1}{2}\rangle S_{b-}|-\frac{1}{2}\rangle) \\
&= (S_{a+}+S_{b+}) \frac{\sigma+1}{\sqrt{2}}|-\frac{1}{2}\rangle|-\frac{1}{2}\rangle \\
&= \frac{\sigma+1}{\sqrt{2}}(S_{a+}|-\frac{1}{2}\rangle|-\frac{1}{2}\rangle+|-\frac{1}{2}\rangle S_{b+}|-\frac{1}{2}\rangle) \\
&= \frac{\sigma+1}{\sqrt{2}}(|+\frac{1}{2}\rangle|-\frac{1}{2}\rangle+|-\frac{1}{2}\rangle|+\frac{1}{2}\rangle)
\end{align}

For this expression to equal zero, we must have $\sigma=-1$. Therefore, the singlet state is:

\begin{equation}
|0,0,\frac{1}{2},\frac{1}{2}\rangle=\frac{1}{\sqrt{2}}(|\frac{1}{2},-\frac{1}{2}\rangle|\frac{1}{2},+\frac{1}{2}\rangle-|\frac{1}{2},+\frac{1}{2}\rangle|\frac{1}{2},-\frac{1}{2}\rangle)
\end{equation}

This state is antisymmetric under particle exchange:

\begin{equation}
|0,0,j_a,j_b\rangle=-|0,0,j_b,j_a\rangle
\end{equation}

For the case $j=1$, we have three possible values: $m_j=-1,0,1$. For $m_j=\pm 1$, there is only one possible configuration:

\[
\left\{\begin{array}{l}
|1,1,j_a,j_b\rangle=|j_a,\frac{1}{2}\rangle|j_b,\frac{1}{2}\rangle=|1,1\rangle  \\
|1,-1,j_a,j_b\rangle=|j_a,-\frac{1}{2}\rangle|j_b,-\frac{1}{2}\rangle=|1,-1\rangle
\end{array}\right.
\]

For $m_j=0$, we again have two possibilities: $m_a=\frac{1}{2}, m_b=-\frac{1}{2}$ or $m_a=-\frac{1}{2}, m_b=\frac{1}{2}$. The state has the form:

\begin{equation}
|1,0,\frac{1}{2},\frac{1}{2}\rangle=\frac{1}{\sqrt{2}}(|j_a,+\frac{1}{2}\rangle|j_b,-\frac{1}{2}\rangle+\nu|j_a,-\frac{1}{2}\rangle|j_b,+\frac{1}{2}\rangle)
\end{equation}

To determine $\nu$, we can generate this state by applying $J_+$ to the $m_j=-1$ state:

\begin{align}
\sqrt{2}|1,0,+\frac{1}{2},+\frac{1}{2}\rangle &= J_{+}|1,-1,+\frac{1}{2},+\frac{1}{2}\rangle=(S_{a+}+S_{b+})|j_a,-\frac{1}{2}\rangle|j_b,-\frac{1}{2}\rangle \\
&= S_{a+}|j_a,-\frac{1}{2}\rangle|j_b,-\frac{1}{2}\rangle+|j_a,-\frac{1}{2}\rangle S_{b+}|j_b,-\frac{1}{2}\rangle \\
&= |j_a,+\frac{1}{2}\rangle|j_b,-\frac{1}{2}\rangle+|j_a,-\frac{1}{2}\rangle|j_b,+\frac{1}{2}\rangle
\end{align}

This calculation shows that $\nu=1$. We can verify this by applying $J_+$ to the $m_j=0$ state:

\begin{align}
\sqrt{2}|1,1,+\frac{1}{2},+\frac{1}{2}\rangle &= J_{+}|1,0,+\frac{1}{2},+\frac{1}{2}\rangle=(S_{a+}+S_{b+})\frac{1}{\sqrt{2}}(|j_a,+\frac{1}{2}\rangle|j_b,-\frac{1}{2}\rangle+|j_a,-\frac{1}{2}\rangle|j_b,+\frac{1}{2}\rangle) \\
&= \frac{1}{\sqrt{2}}(S_{a+}|j_a,+\frac{1}{2}\rangle|j_b,-\frac{1}{2}\rangle+S_{a+}|j_a,-\frac{1}{2}\rangle|j_b,+\frac{1}{2}\rangle \\
&\quad +|j_a,+\frac{1}{2}\rangle S_{b+}|j_b,-\frac{1}{2}\rangle+|j_a,-\frac{1}{2}\rangle S_{b+}|j_b,+\frac{1}{2}\rangle) \\
&= \frac{1}{\sqrt{2}}(|j_a,+\frac{1}{2}\rangle|j_b,+\frac{1}{2}\rangle+|j_a,+\frac{1}{2}\rangle|j_b,+\frac{1}{2}\rangle) \\
&= \sqrt{2}|j_a,+\frac{1}{2}\rangle|j_b,+\frac{1}{2}\rangle
\end{align}

This confirms our choice of $\nu=1$ and reproduces the $m_j=1$ state correctly.

The complete set of triplet states is:

\[
\left\{\begin{array}{l}
|1,1,j_a,j_b\rangle=|j_a,\frac{1}{2}\rangle|j_b,\frac{1}{2}\rangle=|1,1\rangle  \\
|1,0,j_a,j_b\rangle=\frac{1}{\sqrt{2}}(|j_a,+\frac{1}{2}\rangle|j_b,-\frac{1}{2}\rangle+|j_a,-\frac{1}{2}\rangle|j_b,+\frac{1}{2}\rangle) \\
|1,-1,j_a,j_b\rangle=|j_a,-\frac{1}{2}\rangle|j_b,-\frac{1}{2}\rangle=|1,-1\rangle
\end{array}\right.
\]


\section{Matrix formulation for angular momentum operators}
The spectrum-generating algebra approach naturally leads to matrix representations of angular momentum operators. The dimension of these matrices depends on the representation index $j$, requiring $(2j+1)\times(2j+1)$ matrices. For spin-$\frac{1}{2}$ systems, this yields $2\times2$ matrices (the familiar Pauli matrices).

For the case $j=1$, we need $3\times3$ matrices to represent the operators, corresponding to the three basis states with $m_j=-1,0,1$. These basis vectors can be written as:

\[
|1,+1\rangle=\left[\begin{array}{l}
1  \\
0 \\
0
\end{array}\right], \quad|1,0\rangle=\left[\begin{array}{l}
0 \\
1 \\
0
\end{array}\right], \quad|1,-1\rangle=\left[\begin{array}{l}
0 \\
0 \\
1
\end{array}\right]
\]

The matrix representation of $J_3$ in this basis is diagonal:

\[
J_{3}=\hbar\left[\begin{array}{ccc}
1 & 0 & 0  \\
0 & 0 & 0 \\
0 & 0 & -1
\end{array}\right]
\]

We can verify this acts correctly on the basis states, for example:

\[
J_{3}|1,-1\rangle=\hbar\left[\begin{array}{ccc}
1 & 0 & 0  \\
0 & 0 & 0 \\
0 & 0 & -1
\end{array}\right]\left[\begin{array}{l}
0 \\
0 \\
1
\end{array}\right]=-\hbar\left[\begin{array}{l}
0 \\
0 \\
1
\end{array}\right]
\]

To construct the ladder operators $J_+$ and $J_-$, we use the general formulas:

\begin{align}
J_{+}|j,m_j\rangle &= \hbar\sqrt{j(j+1)-m_j(m_j+1)}|j,m_j+1\rangle  \\
J_{-}|j,m_j\rangle &= \hbar\sqrt{j(j+1)-m_j(m_j-1)}|j,m_j-1\rangle
\end{align}

Applying these to our $j=1$ basis states:

\begin{align}
J_{+}|1,+1\rangle &= \hbar\sqrt{1(1+1)-1(1+1)}|1,+2\rangle=0 \\
J_{+}|1,0\rangle &= \hbar\sqrt{1(1+1)-0(0+1)}|1,+1\rangle=\hbar\sqrt{2}|1,+1\rangle \\
J_{+}|1,-1\rangle &= \hbar\sqrt{1(1+1)-(-1)(0)}|1,0\rangle=\hbar\sqrt{2}|1,0\rangle \\
J_{-}|1,+1\rangle &= \hbar\sqrt{1(1+1)-1(0)}|1,0\rangle=\hbar\sqrt{2}|1,0\rangle  \\
J_{-}|1,0\rangle &= \hbar\sqrt{1(1+1)-0(-1)}|1,-1\rangle=\hbar\sqrt{2}|1,-1\rangle \\
J_{-}|1,-1\rangle &= \hbar\sqrt{1(1+1)-(-1)(-2)}|1,-2\rangle=0
\end{align}

To determine the matrix elements, we start with a general form for $J_+$:

\begin{gather}
J_{+}=\hbar\left[\begin{array}{lll}
a & b & c \\
d & e & f \\
g & h & i
\end{array}\right]
\end{gather}

Then we apply this to each basis vector. For example:

\begin{gather}
J_{+}|1,+1\rangle=\hbar\left[\begin{array}{lll}
a & b & c \\
d & e & f \\
g & h & i
\end{array}\right]\left[\begin{array}{l}
1 \\
0 \\
0
\end{array}\right]=\hbar\left[\begin{array}{l}
a \\
d \\
g
\end{array}\right]
\end{gather}


From these conditions, we determine that $a=0, d=0, g=0$. Completing the calculations for $J_+$ and $J_-$:

\[
J_{+}=\hbar \sqrt{2}\left[\begin{array}{lll}
0 & 1 & 0  \\
0 & 0 & 1 \\
0 & 0 & 0
\end{array}\right], \quad J_{-}=\hbar \sqrt{2}\left[\begin{array}{lll}
0 & 0 & 0 \\
1 & 0 & 0 \\
0 & 1 & 0
\end{array}\right]
\]

This reveals the origin of the $\sqrt{2}$ factor encountered in our triplet state calculations. For the spin-$\frac{1}{2}$ case, the corresponding factor would be:

\begin{equation}
\hbar \sqrt{\frac{1}{2}(\frac{1}{2}+1)-\frac{1}{2}(\frac{1}{2}-1)}=\hbar
\end{equation}

This explains the structure of the spin raising and lowering operators.
